\documentclass[a4paper,12pt]{article}
\usepackage[T2A]{fontenc}
\usepackage[utf8]{inputenc}
\usepackage{cmap}
\usepackage[english, russian]{babel}
\usepackage{wrapfig}
\usepackage[backend=biber,style=gost-numeric]{biblatex} % стиль ГОСТ
\addbibresource{literatura.bib}
\usepackage{misccorr} % в заголовках появляется точка, но при ссылке на них ее нет
\usepackage{amssymb,amsfonts,amsmath,amsthm}  
\usepackage{indentfirst}
\usepackage[usenames,dvipsnames]{color} 
\usepackage[unicode,hidelinks]{hyperref}
% \hypersetup{%
%     pdfborder = {0 0 0}
\usepackage{multirow}
\usepackage{makecell,multirow} 
\usepackage{ulem}
%\usepackage{hyperref}
%\usepackage{graphicx,wrapfig}
%\renewcommand\epsilon{\varepsilon}
% \renewcommand\epsilon{\varepsilon}
%\graphicspath{{img/}}
\usepackage{geometry}
\geometry{left=1cm,right=1cm,top=2cm,bottom=2cm,bindingoffset=0cm,headheight=15pt}
\usepackage{fancyhdr} 
\linespread{1.2} 
\frenchspacing 
\renewcommand{\labelenumii}{\theenumii)} 
% \usepackage{caption}
%%%%%%%%%%%%%%%%%%%%%%%%%%%%%%%%%%%%%%%%%%%%%%%%%%%%%%%%%%%%%%%%%%%%%%%%%%%%%%%
%%%%%%%%%%%%%%%%%%%%%%%%%%%%%%%%%%%%%%%%%%%%%%%%%%%%%%%%%%%%%%%%%%%%%%%%%%%%%%%

\def\labauthor{Сарафанов И.Г.}
\def\labauthors{\labauthor}
\def\labnumber{127}
\def\labtheme{Таблица производных и интегралов}

%%%%%%%%%%%%%%%%%%%%%%%%%%%%%%%%%%%%%%%%%%%%%%%%%%%%%%%%%%%%%%%%%%%%%%%%%%%%%%%
	%применим колонтитул к стилю страницы
\pagestyle{fancy} 
	%очистим "шапку" страницы
\fancyhead{\LaTeX} 
	%слева сверху на четных и справа на нечетных
\fancyhead[L]{\labauthors} 
	%справа сверху на четных и слева на нечетных
\fancyhead[R]{Иродов \textnumero 1.185} 
	%очистим "подвал" страницы
\fancyfoot{} 
	% номер страницы в нижнем колинтуле в центре
\fancyfoot[C]{\thepage} 
%%%%%%%%%%%%%%%%%%%%%%%%%%%%%%%%%%%%%%%%%%%%%%%%%%%%%%%%%%%%%%%%%%%%%%%%%%%%%%%
\usepackage{float}
\usepackage[mode=buildnew]{standalone}
\usepackage{tikz} 
% \usepackage{subcaption}
\usepackage{tikz,csvsimple,physics}
\makeatother
%\usepackage{booktabs}
\usepackage{pgfplots, pgfplotstable}
\usepackage[outline]{contour}
\usepackage{tocloft}
\renewcommand{\cftsecleader}{\cftdotfill{\cftdotsep}}

\begin{document}
\section*{Задача 1.186}

\textbf{Условие:} Найти минимальную скорость нейтрона массы $m$, чтобы при столкновении с покоящимся ядром массы $M$ увеличить его внутреннюю энергию на $\Delta E$.

\bigskip

\textbf{Решение:}

\medskip
\textbf{Шаг 1. Определяем законы, которые используем.}

Так как процесс неупругий (ядро получает внутреннюю энергию $\Delta E$):

\begin{itemize}
    \item Сохраняется импульс системы: 
    \[
        \vec p_\text{до} = \vec p_\text{после}
    \]
    \item Полная энергия сохраняется, но кинетическая энергия частично может перейти во внутреннюю энергию ядра:
    \[
        W_\text{полная} = W_\text{мех} + W_\text{внутр}
    \]
\end{itemize}

\medskip
\textbf{Шаг 2. Кинетическая энергия системы до столкновения.}

Нейтрон движется со скоростью $v$, ядро покоится:

\[
T_\text{до} = \frac{1}{2} m v^2 + \frac{1}{2} M \cdot 0^2 = \frac{1}{2} m v^2
\]

\medskip
\textbf{Шаг 3. Разложение энергии через центр масс (теорема Кёнига).}

Скорость центра масс системы:

\[
V_\text{cm} = \frac{m v + M \cdot 0}{m + M} = \frac{m}{m + M} v
\]

Скорость нейтрона относительно ядра:

\[
v_\text{rel} = v - 0 = v
\]

Полная кинетическая энергия раскладывается на:

\[
T = \underbrace{\frac{1}{2}(m+M)V_\text{cm}^2}_{T_\text{мех}} + \underbrace{\frac{1}{2} \mu v_\text{rel}^2}_{T_\text{отн}}
\]

где приведённая масса:

\[
\mu = \frac{m M}{m + M}
\]

\medskip
\textbf{Шаг 4. Вычисляем каждую часть.}

\[
T_\text{мех} = \frac{1}{2} (m+M) V_\text{cm}^2 = \frac{1}{2} (m+M) \left(\frac{m}{m+M} v\right)^2 = \frac{m^2}{2 (m+M)} v^2
\]

\[
T_\text{отн} = \frac{1}{2} \mu v_\text{rel}^2 = \frac{1}{2} \frac{m M}{m+M} v^2
\]

\medskip
\textbf{Шаг 5. Связь относительной энергии и внутренней энергии.}

Энергия относительного движения $T_\text{отн}$ может полностью перейти во внутреннюю энергию ядра. Минимальная скорость достигается, когда после столкновения частицы движутся вместе (полностью неупругий удар):

\[
T_\text{отн} = \Delta E
\]

Подставляем выражение для $T_\text{отн}$:

\[
\frac{1}{2} \frac{m M}{m+M} v^2 = \Delta E
\]

---

\textbf{Шаг 6. Находим минимальную скорость нейтрона.}

\[
v^2 = \frac{2 \Delta E (m+M)}{m M} 
\]

\[
\boxed{v_\text{min} = \sqrt{\frac{2 \Delta E (m+M)}{m M}}}
\]

---

\textbf{Шаг 7. Физический смысл решения.}

\begin{itemize}
    \item Движение центра масс системы $V_\text{cm}$ не может быть использовано для возбуждения ядра, оно остаётся механическим. 
    \item Внутренняя энергия ядра $\Delta E$ может быть получена только за счёт энергии относительного движения $T_\text{отн} = \frac{1}{2} \mu v_\text{rel}^2$. 
    \item Приведённая масса $\mu$ показывает, какая часть кинетической энергии доступна для преобразования во внутреннюю энергию.
    \item Если ядро сильно тяжелее нейтрона ($M \gg m$), $\mu \approx m$ и почти вся энергия нейтрона может уйти во внутреннюю энергию.
\end{itemize}

\bigskip
\textbf{Ответ:}

\[
\displaystyle v_\text{min} = \sqrt{\frac{2 \Delta E (m+M)}{m M}}
\]
\end{document}